\section{主要研究方法}\label{sec:methods}

\subsection{圆法（Circle Method）}

圆法由哈代和李特尔伍德于1920年代发展\cite{hardy1923some}，是研究哥德巴赫猜想最重要的方法之一。

\subsubsection{基本思想}

设 $r(n)$ 表示偶数 $n$ 表示为两个素数之和的方法数，则
\begin{equation}\label{eq:circle_method}
    r(n) = \int_0^1 S(\alpha)^2 e(-n\alpha) \, d\alpha
\end{equation}
其中
\begin{equation}
    S(\alpha) = \sum_{p \leq n} (\log p) e(p\alpha)
\end{equation}
是素数的指数和，$e(x) = e^{2\pi i x}$。

\subsubsection{优弧与劣弧}

将积分区间 $[0,1]$ 分为两部分：
\begin{itemize}
    \item \textbf{优弧}（Major arcs）：$\alpha$ 接近有理数 $a/q$（$q$ 较小）的区域
    \item \textbf{劣弧}（Minor arcs）：其余区域
\end{itemize}

优弧上的积分贡献主要项，劣弧上的积分需要证明是低阶项。

\subsection{筛法（Sieve Methods）}

筛法是研究哥德巴赫猜想的另一重要工具。

\subsubsection{埃拉托斯特尼筛法}

经典的埃拉托斯特尼筛法通过逐步剔除合数来筛选素数，其思想是现代筛法的源头。素数定理给出
\begin{equation}
    \pi(x) \sim \frac{x}{\log x}, \qquad x \to \infty,
\end{equation}
即 $x$ 附近的素数密度约为 $1/\log x$，这是估计筛法“筛剩”比例的基础\cite{davenport2000multiplicative}。

\subsubsection{布伦筛法}

布伦筛法（1920）是第一个系统化的组合筛法\cite{brun1920sieve}，系统阐述见\cite{halberstam1974sieve}。其核心不等式为：
\begin{theorem}[布伦筛法]
    设 $\mathcal{A}$ 为有限整数集合，$\mathcal{P}$ 为素数集合，则存在筛法函数 $S(\mathcal{A}, \mathcal{P}, z)$ 满足：
    \begin{equation}
        S(\mathcal{A}, \mathcal{P}, z) = \sum_{\substack{a \in \mathcal{A} \\ (a, P(z))=1}} 1
    \end{equation}
    其中 $P(z) = \prod_{\substack{p < z \\ p \in \mathcal{P}}} p$。
\end{theorem}

\subsubsection{塞尔伯格筛法}

塞尔伯格筛法（1947）通过选择最优筛法权重改进了布伦筛法\cite{selberg1949sieve}。其上界为：
\begin{equation}
    S(\mathcal{A}, \mathcal{P}, z) \leq \frac{X}{V(z)} + O\left(\sum_{\substack{d_1, d_2 < z \\ d_1, d_2 | P(z)}} \lambda_{d_1} \lambda_{d_2} r_{[d_1, d_2]}\right)
\end{equation}

\subsection{指数和估计}

指数和估计在哥德巴赫猜想研究中起着关键作用。

\subsubsection{华罗庚引理}

华罗庚引理给出了 $k$ 次幂指数和的一个基本上界：
\begin{equation}
    \int_0^1 \left| \sum_{n=1}^N e(\alpha n^k) \right|^{2^k} \, d\alpha \ll N^{2^k - k + \epsilon},
\end{equation}
对任意 $\epsilon > 0$ 成立。它是圆法中控制劣弧的关键输入之一\cite{iwaniec2004analytic}。

\subsubsection{外尔不等式与几何级数界}

对线性相位和，有精确的几何级数估计：
\begin{equation}
    \left| \sum_{n=1}^N e(\alpha n) \right| \le \min\!\left(N,\ \frac{1}{2\|\alpha\|}\right),
\end{equation}
其中 $\|\alpha\|$ 表示 $\alpha$ 到最近整数的距离。对于多项式相位 $\alpha_k n^k + \cdots + \alpha_1 n$，外尔不等式给出次凸型上界：当 $\alpha_k = a/q$（互素）时，
\begin{equation}
    \left| \sum_{n=1}^N e(\alpha_k n^k + \cdots + \alpha_1 n) \right| \ll N^{1+\epsilon} \left( \frac{1}{q} + \frac{1}{N} + \frac{q}{N^k} \right)^{1/2^{k-1}}.
\end{equation}

\subsection{大筛法}

大筛法是解析数论中的重要工具，由林尼克和邦别里发展\cite{iwaniec2004analytic}。

\begin{theorem}[大筛法不等式]
    设 $a_1, \ldots, a_N$ 为复数，则
    \begin{equation}
        \sum_{q \leq Q} \sum_{\substack{a=1 \\ (a,q)=1}}^q \left| \sum_{n=1}^N a_n e\left(\frac{an}{q}\right) \right|^2 \leq (N + Q^2) \sum_{n=1}^N |a_n|^2
    \end{equation}
\end{theorem}
